%% 39장 특수함수 I
\chapter{특수함수 I: 감마와 오차함수}
\label{ch:special-1}

\abstract{통계와 물리의 공식에는 기본 함수만으로는 나타낼 수 없는 함수들이 되풀이해 나타난다. 이 장과 다음 두 장에서는 그런 특수함수를 binary64에서 상대오차 $10^{-14}$ 안으로 계산하는 방법을 다룬다. 이 장의 주인공은 감마 함수와 오차함수이다. 감마 함수의 함수 방정식과 반사 공식, Weierstrass 곱에서 나오는 $\log\Gamma$의 거듭제곱 급수를 증명하고, 점근급수인 Stirling 급수가 항을 늘리면 언젠가 발산한다는 것을 정확한 계산으로 보인다. Lanczos 근사가 넓은 구간을 한 식으로 덮는 방식, $\log\Gamma$가 $0$이 되는 $1$과 $2$ 근처에서 상대 정확도를 지키는 급수, 베타 함수와 디감마 함수를 본다. 오차함수에서는 보수 함수를 따로 계산해야 하는 이유와, $e^{-x^2}$에서 $x^2$의 반올림이 $2x^2$배로 커지는 문제를 융합 곱셈--덧셈으로 푸는 방법을 본다.}

\section{감마 함수}
\label{sec:gamma}

\subsection{정의와 함수 방정식}

$\operatorname{Re}z>0$에서 \term{감마 함수}{gamma function}는 Euler의 적분
\[
\Gamma(z)=\int_0^\infty t^{z-1}e^{-t}\,dt
\]
로 정의된다\cite{davis1959gamma}. 부분적분을 하면 경계항이 사라지고 $\Gamma(z+1)=z\Gamma(z)$를 얻는다. $\Gamma(1)=1$이므로 양의 정수에서는 $\Gamma(n+1)=n!$이다. 이 함수 방정식을 거꾸로 쓰면 $\Gamma(z)=\Gamma(z+1)/z$로 왼쪽 반평면으로 해석적 연장을 할 수 있고, $0,-1,-2,\dots$에 단순 극점이 생긴다. 실수에서 계산할 때도 이 점화식이 기본 도구이다. 큰 $x$에서 잘 맞는 근사를 작은 $x$로 끌어내리거나 그 반대를 할 수 있기 때문이다.

\cref{eq:reflection}의 \term{반사 공식}{reflection formula}은 음의 실수를 양의 실수로 옮긴다.
\begin{equation}
\label{eq:reflection}
\Gamma(z)\Gamma(1-z)=\frac\pi{\sin\pi z}
\end{equation}
이 공식은 아래의 Weierstrass 곱과 사인 함수의 곱 표현 $\sin\pi z=\pi z\prod_{k\ge1}(1-z^2/k^2)$에서 바로 나온다\cite{whittaker1927}. 계산할 때 주의할 점은 $\sin\pi x$이다. $\pi x$를 binary64로 만들면 $x$가 클 때 그 반올림이 $\sin$의 값을 완전히 바꾼다. 그래서 인자를 $x$의 정수 부분에서 정확히 줄인 뒤 $\sin(\pi r)$, $|r|\le1/2$을 계산하는 \code{sinpi} 같은 함수를 써야 한다. 8장에서 본 대로 $x-\operatorname{round}(x)$는 정확하다.

실수 양의 반직선 위에서 감마 함수는 다음 성질로 완전히 특징지어진다. $\Gamma(1)=1$, $\Gamma(x+1)=x\Gamma(x)$이고 $\log\Gamma$가 볼록인 함수는 감마 함수뿐이다. 이것이 \term{Bohr--Mollerup 정리}{Bohr--Mollerup theorem}이다\cite{artin1964}.

\subsection{Weierstrass 곱과 $\log\Gamma$의 급수}

감마 함수는 해석 함수 $1/\Gamma$의 곱 표현
\begin{equation}
\label{eq:weierstrass}
\frac1{\Gamma(z)}=ze^{\gamma z}\prod_{k=1}^\infty\Bigl(1+\frac zk\Bigr)e^{-z/k}
\end{equation}
을 가진다. 여기서 $\gamma=0.5772\dots$는 Euler--Mascheroni 상수이다\cite{whittaker1927}. 이것에서 $1$ 근처의 거듭제곱 급수가 나온다.

\begin{proposition}
\label{prop:loggamma-series}
$|z|<1$에서
\[
\log\Gamma(1+z)=-\gamma z+\sum_{k=2}^\infty\frac{(-1)^k\zeta(k)}kz^k
\]
이다. 따라서 $|e|<1$에서
\[
\log\Gamma(2+e)=(1-\gamma)e+\sum_{k=2}^\infty\frac{(-1)^k\bigl(\zeta(k)-1\bigr)}ke^k
\]
이다.
\end{proposition}

\begin{proof}
$\Gamma(1+z)=z\Gamma(z)$이므로 \eqref{eq:weierstrass}에서 $-\log\Gamma(1+z)=\gamma z+\sum_k\bigl(\log(1+z/k)-z/k\bigr)$이다. $\log(1+w)-w=\sum_{m\ge2}(-1)^{m+1}w^m/m$을 넣고 $k$에 대한 합과 $m$에 대한 합을 바꾸면(절대 수렴한다) $\sum_{m\ge2}(-1)^{m+1}\zeta(m)z^m/m$이 되어 첫 식을 얻는다. 둘째 식은 $\log\Gamma(2+e)=\log(1+e)+\log\Gamma(1+e)$에 $\log(1+e)=\sum_{k\ge1}(-1)^{k+1}e^k/k$를 더하면 나온다.
\end{proof}

이 급수는 계산에서 특별한 쓸모가 있다. $\log\Gamma$는 $x=1$과 $x=2$에서 $0$이므로, 그 근처에서 값은 작고 상대 정확도를 지키기 어렵다. $\log(\Gamma(x))$처럼 $\Gamma$를 먼저 계산하면 $\Gamma(x)\approx1$의 반올림 오차 $u$가 로그를 거쳐 절대오차 $u$로 남는데, 값 자체가 $(1-\gamma)(x-2)$ 정도이므로 상대오차는 $u/|x-2|$로 커진다. $x=2+10^{-12}$에서 $\log(\Gamma(x))$의 상대오차는 $1.2\times10^{-4}$이고, 이 Mac의 C 수학 라이브러리의 \code{lgamma}도 $1.6\times10^{-4}$이다. \cref{prop:loggamma-series}의 둘째 식은 $e$의 배수로 시작하므로 상대 정확도를 그대로 지킨다. 게다가 $\zeta(k)-1\approx2^{-k}$이어서 $|e|\le1/2$이면 항이 빨리 준다.

\section{점근 전개}
\label{sec:stirling}

\subsection{Stirling 급수}

큰 $z$에서는 \term{Stirling 급수}{Stirling series}
\begin{equation}
\label{eq:stirling}
\log\Gamma(z)=\Bigl(z-\frac12\Bigr)\log z-z+\frac12\log2\pi+\sum_{k=1}^{N}\frac{B_{2k}}{2k(2k-1)z^{2k-1}}+R_N(z)
\end{equation}
가 쓰인다. $B_{2k}$는 Bernoulli 수이다. 이 급수는 수렴하지 않는다. $|B_{2k}|$가 $(2k)!/(2\pi)^{2k}$의 두 배쯤으로 자라므로 고정된 $z$에서 항은 언젠가 커진다. 그러나 양의 실수 $z$에서 나머지 $R_N$은 처음 버린 항과 부호가 같고 크기가 그 이하이다\cite{olver1974asymptotics}. 따라서 항이 가장 작아지는 곳, 대략 $N\approx\pi z$에서 멈추면 최선의 정확도를 얻는다. 이런 급수를 3장에서 본 \term{점근급수}{asymptotic series}라 한다.

$z=5$에서 $\log\Gamma(5)=\log24$와 견주어 보자. Python의 \code{decimal}로 $50$자리에서 계산하면 항을 $1,2,4,8,12$개 쓸 때 오차가 $2.2\times10^{-5}$, $2.5\times10^{-7}$, $4.0\times10^{-10}$, $1.8\times10^{-13}$, $4.5\times10^{-15}$로 줄고, $16$개($\approx\pi z$)에서 $1.6\times10^{-15}$로 가장 작으며, $20$개에서 $4.1\times10^{-15}$, $24$개에서 $5.2\times10^{-14}$, $30$개에서 $2.7\times10^{-11}$로 다시 커진다. $z=5$에서는 binary64의 정확도에 이르지 못한다. 그래서 실제 구현은 점화식으로 $z$를 충분히 키운 뒤 몇 개의 항만 쓴다. 점화식으로 키우면 $\Gamma(z)=\Gamma(z+m)/(z(z+1)\cdots(z+m-1))$의 곱 $m$개가 반올림 오차를 더한다.

\subsection{Lanczos 근사}

Lanczos는 Stirling 급수의 앞부분에 유리함수를 곱한 꼴로 오른쪽 반평면 전체를 한 식으로 덮는 근사를 내놓았다\cite{lanczos1964}.
\[
\Gamma(z+1)=\sqrt{2\pi}\Bigl(z+g+\tfrac12\Bigr)^{z+1/2}e^{-(z+g+1/2)}\Bigl(c_0+\sum_{k=1}^{n}\frac{c_k}{z+k}+\epsilon_{g,n}(z)\Bigr)
\]
매개변수 $g$를 고르면 계수 $c_k$가 정해지고, 적당한 $g$에서는 항 몇 개만으로 오차 $\epsilon$이 오른쪽 반평면 전체에서 아주 작다. Pugh는 $g$와 $n$의 선택과 오차를 체계적으로 분석했다\cite{pugh2004}. 널리 퍼진 $g=7$, $n=9$의 계수로 $[0.6,40]$에서 $\Gamma$를 계산해 위의 \code{decimal} 참값과 견주면 상대오차가 $3.8\times10^{-14}$ 이하이다. 계산이 고정된 수의 나눗셈과 지수, 거듭제곱 몇 번이므로 빠르고, 분기가 없어 벡터화하기 쉽다. Boost.Math는 $g\approx6.02$의 $13$항 판을 쓴다\cite{boostmath}.

\subsection{조건수}

정확도를 말하기 전에 문제 자체의 조건수를 보자. 상대 조건수는 $|x\Gamma'(x)/\Gamma(x)|=|x\psi(x)|$이고, 큰 $x$에서 $x\log x$ 정도로 자란다. $x=100$이면 약 $460$이므로 $x$에 $u$만큼의 오차가 있으면 $\Gamma(x)$는 $460u$만큼 틀린다. 반올림된 입력 $x$ 자체는 정확하다고 보므로 라이브러리의 목표는 그 입력의 정확한 $\Gamma$에 가까운 값이지만, 계산 중에 $x+g+1/2$ 같은 중간값을 반올림하면 그 오차가 이 조건수만큼 커진다. 그래서 정확한 구현은 그런 중간값을 두 binary64의 합으로 들고 간다. 극점 근처의 음수 $x$에서도 조건수가 커지는데, 이것은 반사 공식의 $\sin\pi x$가 $0$에 가깝기 때문이다.

\section{베타 함수와 디감마 함수}
\label{sec:beta-digamma}

\term{베타 함수}{beta function} $B(a,b)=\int_0^1t^{a-1}(1-t)^{b-1}dt=\Gamma(a)\Gamma(b)/\Gamma(a+b)$를 감마 함수 셋으로 계산하면 두 가지 문제가 생긴다. 큰 인자에서 분자와 분모가 따로 넘치고, 세 감마 함수의 반올림 오차가 모인다. 로그를 써서 $\log\Gamma$ 셋의 합으로 계산하면 넘치지는 않지만, 큰 값들의 뺄셈에서 9장의 소거가 일어난다. Lanczos 근사를 쓰면 세 식의 지수 인자들을 한 식에 모아 $\bigl((a+g-1/2)/(c+g-1/2)\bigr)^{a-1/2}$ 같은 비율의 거듭제곱으로 정리할 수 있어 두 문제를 함께 피한다.

\term{디감마 함수}{digamma function} $\psi=\Gamma'/\Gamma$는 \eqref{eq:weierstrass}를 로그 미분해 $\psi(z)=-\gamma+\sum_{k\ge0}\bigl(1/(k+1)-1/(k+z)\bigr)$로 쓸 수 있다. 계산은 감마와 같은 틀을 따른다. 음수는 반사 공식을 미분한 $\psi(1-x)-\psi(x)=\pi\cot\pi x$로 옮기고, 작은 양수는 점화식 $\psi(x+1)=\psi(x)+1/x$로 키운 뒤 점근급수
\[
\psi(z)\sim\log z-\frac1{2z}-\sum_{k\ge1}\frac{B_{2k}}{2kz^{2k}}
\]
를 쓴다. 다만 $\psi$는 양의 실수에서 $x_0=1.4616321449683623\dots$에 근이 하나 있어서, 그 근처에서 상대 정확도를 지키려면 근을 중심으로 전개한 급수를 따로 써야 한다\cite{cody1973psi}. $\log\Gamma$가 $1$과 $2$에서 겪은 문제와 같다. 함수의 근 근처에서 상대오차를 지키는 일은 특수함수 구현의 되풀이되는 주제이다.

\section{오차함수}
\label{sec:erf}

\subsection{보수 함수를 따로}

\term{오차함수}{error function} $\operatorname{erf}x=\frac2{\sqrt\pi}\int_0^xe^{-t^2}dt$는 정규분포의 누적분포함수 $\Phi(x)=\frac12\operatorname{erfc}(-x/\sqrt2)$의 바탕이다. 큰 $x$에서 $\operatorname{erf}x$는 $1$에 아주 가까우므로, \term{여오차함수}{complementary error function} $\operatorname{erfc}x=1-\operatorname{erf}x$를 $1$에서 빼서 계산하면 9장의 소거로 상대 정확도를 모두 잃는다. $x=6$이면 $\operatorname{erfc}6\approx2.2\times10^{-17}$은 $1-\operatorname{erf}6$을 binary64로 계산하면 $0$이다. 그래서 $\operatorname{erfc}$는 따로 근사해야 하고, 큰 $x$에서는 $\operatorname{erfc}x=e^{-x^2}\cdot R(x)$ 꼴의 유리 근사를 쓴다. Cody는 세 구간에서 Chebyshev 유리 근사를 만들었고\cite{cody1969erf}, Cephes 라이브러리\cite{moshier1989}의 \code{erf}와 \code{erfc}도 비슷한 세 구간 구조이다. 꼬리의 정규 확률을 다루는 통계 계산에서 이 구분은 결정적이다. 정규분포의 분위수 함수 $\Phi^{-1}$은 Wichura의 AS 241이 가운데 구간과 두 꼬리 구간의 유리 근사로 binary64의 정확도를 낸다\cite{wichura1988}.

\term{척도 여오차함수}{scaled complementary error function} $\operatorname{erfcx}x=e^{x^2}\operatorname{erfc}x$도 자주 필요하다. $\operatorname{erfc}$는 $x$가 $27$ 근처를 넘으면 아래로 넘치지만 $\operatorname{erfcx}$는 $1/(x\sqrt\pi)$처럼 천천히 줄므로, 꼬리 확률의 로그나 비율을 넘침 없이 계산할 수 있다.

\subsection{$e^{-x^2}$의 반올림}

$\operatorname{erfc}x=e^{-x^2}R(x)$에서 $x^2$을 binary64로 계산하면 상대오차 $u$, 곧 절대오차 $ux^2$가 생긴다. 이것이 지수 함수를 지나면 $e^{-x^2}$의 상대오차 $ux^2$로 커진다. $e^{-y}$의 상대 조건수가 $|y|$이기 때문이다. $x=26$이면 $ux^2\approx7.5\times10^{-14}$이어서 목표인 $10^{-14}$을 넘는다. 8장의 융합 곱셈--덧셈은 이 문제를 깔끔하게 푼다. $h=\operatorname{fl}(x\cdot x)$와 $l=\operatorname{fma}(x,x,-h)$로 $x^2=h+l$을 정확히 나누면 $e^{-x^2}=e^{-h}e^{-l}\approx e^{-h}(1-l)$이고, $|l|\le u h$이므로 $e^{-l}$의 1차 근사로 충분하다. 직접 재어 보면 $e^{-x^2}$의 상대오차가 $x=10.3$에서 $3.0\times10^{-15}$에서 $8.5\times10^{-17}$로, $x=26.2$에서 $2.2\times10^{-14}$에서 $9.0\times10^{-17}$로 준다. Cephes는 FMA가 흔하지 않던 시절 $x$를 정수 배의 $1/128$과 나머지로 나누어 같은 효과를 얻었다.

\section{stardust의 감마와 오차함수}
\label{sec:special1-stardust}

\begin{instardust}[\code{docs/design/special.md} §2.3, §5.2]
stardust의 \code{sd.special}은 함수마다 보장 영역을 정하고 그 안에서 Arb 참값 대비 상대오차 $10^{-14}$ 이하를 계약으로 삼는다. \code{gamma}는 Boost의 $13$항 Lanczos 식을 쓰되 $z+g-1/2$를 두 binary64의 합으로 들고 다니고, 큰 $z$에서는 거듭제곱을 반씩 나누어 넘치지 않게 한다. 음수는 반사 공식으로 옮기는데 $\sin\pi x$는 인자를 정확히 줄이는 Apple의 \code{\_\_sinpi}로 계산한다. \code{gammaln}은 $2$ 둘레에서 \cref{prop:loggamma-series}의 $\zeta(k)-1$ 계수로 $28$항까지 더하고, 그 밖에서는 Stirling 근사를 쓴다. $1$ 근처는 $\log\Gamma(2+e)-\log(1+e)$로 옮긴다. \code{beta}는 세 Lanczos 식을 한 식에 모은 비율로 계산한다.
\end{instardust}

\begin{instardust}[\code{docs/design/special.md} §5.1]
\code{erf}와 \code{erfc}는 Cephes의 세 구간 유리식을 쓰되, $e^{-x^2}$는 위의 FMA 분해로 계산한다. \code{erfcx}는 $x\ge8$에서 Cephes의 유리식 대신 점근급수를 쓰는데, 설계 문서에 따르면 그 유리식은 $x=26.64$까지 맞춘 식이어서 $x\approx10^4$에서 $21$ulp가 틀렸다. \code{ndtri}는 Wichura의 AS 241이고, \code{erfinv}는 그 가운데 구간의 식을 $q=y/2$로 바로 불러 작은 $y$에서 $(1+y)/2$의 소거를 피한다. \code{log\_ndtr}는 음수 $x$에서 $-x^2/2+\log(\operatorname{erfcx}(-x/\sqrt2)/2)$로 계산해, 아래로 넘치지 않으면서 두 항의 부호가 같아 소거도 없다.
\end{instardust}

\section{맺음}

특수함수를 정확히 계산하는 일은 하나의 공식을 찾는 일이 아니라, 정의역을 나누어 구간마다 알맞은 표현을 고르고 그 사이를 정확한 변환으로 잇는 일이다. 점화식과 반사 공식은 정의역을 옮기고, 거듭제곱 급수와 점근급수와 유리 근사는 각자의 구간을 맡는다. 그리고 함수가 $0$이 되는 곳, 값이 $1$에 붙는 곳, 지수가 반올림을 키우는 곳마다 상대 정확도를 지키는 특별한 손질이 필요하다. 다음 장에서는 이 도구들을 불완전 감마 함수와 불완전 베타 함수로 넓히고, 그 계산에 등장하는 연분수와 균등 점근 전개를 본다.

특수함수의 이론은 Whittaker와 Watson의 고전\cite{whittaker1927}과 Temme의 책\cite{temme1996}이, 공식은 NIST의 DLMF\cite{dlmf}가 표준이다.

\chapterreferences
